\documentclass[../../main.tex]{subfiles}
\begin{document}

{\bf [解]}

与えられた和
\begin{align}
  S = \sum_{k=1}^{100} |n-k| \label{eq:1}
\end{align}
の最小値を考える．絶対値を外すため，$n$の値に応じて場合分けする．

\subparagraph{$n \le 1$ の時}

\cref{eq:1}の絶対値を外して
\begin{align}
  S
  = \sum_{k=1}^{100} (k-n)
  = \dfrac{1}{2} \cdot 100 \cdot 101 - 100 n
\end{align}
だからこれは$n$について単調減少であり，$n=1$で最小値$\min S = 4950$をとる．

\subparagraph{$n \ge 100$ の時}

\cref{eq:1}の絶対値を外して
\begin{align}
  S = \sum_{k=1}^{100} (n-k) = 100 n - \dfrac{1}{2} \cdot 100 \cdot 101
\end{align}
だからこれは $n$ について単調増加であり，$n=100$で最小値 $\min S = 4950$をとる．

\subparagraph{$1 \le n \le 99$の時}

\cref{eq:1}の絶対値を外して
\begin{align}
  S
  &= \sum_{k=1}^n (n-k) + \sum_{k=n+1}^{100} (k-n) \\
  &= n^2 - \dfrac{1}{2} n(n+1) + \dfrac{1}{2}(101-n)(100-n) \\
  &= n^2 - \dfrac{1}{2} n - \dfrac{1}{2} \cdot 201 n + 50 \cdot 101 \\
  &= n^2 - 101 n + 50 \cdot 101 \\
  &= \left(n - \dfrac{101}{2}\right)^2 + 50 \cdot 101 - \dfrac{101^2}{4}
\end{align}
である．これは自然数の$n$に対しては$n = 50, 51$ の時に最小値 $\min S = 2500$ をとる．
\casesend

以上の場合分けから，$S$が最小となるのは$n=50,51$の時の$\min S = 2500$である．

\newpage

\end{document}
